\section[Conclusão]{Conclusão}

Esta sequência de sprints foi marcada por uma mistura de desafios intensos, aprendizado profundo e muita colaboração. Desde o início, com a retrospectiva da Sprint 0, nossa equipe demonstrou resiliência e adaptabilidade, especialmente ao nos dividirmos em subgrupos focados no front-end, back-end e blockchain. Assumi o papel no grupo de back-end, que era de particular interesse para mim.

A primeira sprint trouxe um foco significativo na modelagem do banco de dados, o que revelou desafios devido às especificações voláteis das User Stories. No entanto, conseguimos definir uma estrutura robusta, que incluiu as entidades principais de Empresa, Usuário, Certificado e Permissão. Além disso, assumi tarefas relacionadas a métodos HTTP, o que me permitiu me familiarizar com a aplicação back-end e suas APIs.

Apesar de alguns problemas técnicos com Merge Requests, a colaboração com colegas de AGES foi vital para resolver esses problemas. A experiência de tornar operações assíncronas na blockchain foi particularmente valiosa, monstrando a importância da eficiência e responsividade na aplicação.

Ao longo das sprints, enfrentamos grandes problemas, como a saída de dois membros da equipe, o que aumentou a responsabilidade sobre mim e outro colega no desenvolvimento do banco de dados. Mesmo assim, conseguimos manter a produtividade e a qualidade do trabalho, graças ao apoio mútuo e à dedicação de todos os membros da equipe. Além disso, as enchentes que varreram nosso estado complicaram muito o andamento do projeto, causando interrupções nas aulas e afetando o moral e a saúde de muitos integrantes. Apesar dessas adversidades, conseguimos retomar o ritmo e ajustar nossas estratégias para garantir que o projeto continuasse a progredir.

Foi claro que cada desafio enfrentado foi uma oportunidade de crescimento. A experiência adquirida ao longo deste projeto, desde a modelagem do banco de dados até a implementação de operações assíncronas na blockchain, foi muito valiosa. As lições aprendidas e as habilidades desenvolvidas não só contribuíram para o sucesso deste projeto, mas também prepararam cada um de nós para enfrentar futuros desafios com maior confiança e competência.

Em conclusão, essas sprints não apenas resultaram em um projeto bem-sucedido, mas também fortaleceram nossa coesão como equipe e aprimoraram nossas habilidades técnicas. Estou extremamente grato por essa experiência enriquecedora e ansioso para aplicar esses conhecimentos em projetos futuros.